\section{Conclusion}

This paper re-evaluates attention-based hallucination methods through the lens of evidence verification. For detection, we show that generation position can explain much of the apparent success of attention-based scores, and that position-controlled evaluation sharply separates attention-mass scores from representation and uncertainty baselines. For mitigation, we show that attention interventions can alter output behavior without reliably reducing hallucinations: ClearSight shifts POPE answers toward ``yes,'' PAI-attention-only changes attention without changing decisions, and Visual Attention Sink increases hallucinated object mentions in captions.

The central lesson is simple but important: attention may reveal where a model looks, but it does not by itself establish whether the model verifies what it says. The critical distinction is target-discriminative evidence. A model can attend to a real, semantically associated region and still make a false object claim; amplifying that attention can strengthen plausibility rather than verification. Future hallucination methods should treat attention as a candidate feature to be controlled, selected, or learned from, not as a direct certificate of grounding. A detector or mitigation method should earn its claim by showing that its signal remains selective under position, category, answer-prior, and semantic-neighbor controls.
